\documentclass{article}
\usepackage[margin=1in]{geometry}
\usepackage{amsmath}
\usepackage{amssymb}
\usepackage{booktabs}
\usepackage{array}
\usepackage{hyperref}
\newcommand{\score}[2]{#1{\scriptsize$\pm$#2}}

\title{Encoder-Level Lipschitz Regularization in STEM-GNN: Preliminary Results}
\author{Internal Research Memo \\ Working draft for discussion}
\date{\today}

\begin{document}
\maketitle

\begin{abstract}
STEM-GNN's finetune stage previously regularized only the linear
classification head's Lipschitz constant (via a Frobenius-norm penalty on
its weight matrix, \texttt{decoder\_jac\_coeff}), while the GNN encoder
beneath it remains unfrozen and unconstrained during finetuning. This memo
reports a preliminary comparison of that head-only penalty against a new
matching penalty applied to the encoder backbone's own weight matrices
(\texttt{encoder\_lip\_coeff}), added directly to \texttt{model/encoder.py}.
Results are from a reduced-budget run on Cora and should be read as a first
look, not a final claim.
\end{abstract}

\section{Motivation}

The finetune objective already included a Jacobian-style penalty on the
decoder,
\[
  \mathcal{L}_{\text{head}} = \lambda_{\text{dec}} \lVert W_{\text{dec}} \rVert_F^2,
\]
intended to bound the sensitivity of the final linear map to its input.
However, during finetuning the encoder itself is not frozen (only the
vector-quantizer is, by default) and carries no such constraint. An
encoder that is free to amplify its input by an arbitrary factor can pass a
perturbation through to the head no matter how tightly the head's own
weight norm is controlled, so a head-only penalty bounds only the last link
in the chain. We added an analogous penalty over the encoder's backbone
weight matrices,
\[
  \mathcal{L}_{\text{enc}} = \lambda_{\text{enc}} \sum_{\ell}
  \lVert W_\ell \rVert_F^2,
\]
summed over every weight matrix in the backbone's message-passing layers
(SAGE/GAT/GCN/GIN, including per-expert weights when MoE layers are used),
excluding biases and GAT's attention vectors. Both penalties are gated by
independent coefficients and default to $0$ (opt-in, no behavior change
unless requested).

\section{Experimental Setup}

\paragraph{Honesty note on scope.} This is a reduced-budget, single-dataset
sanity check, not the project's standard evaluation protocol. It was run on
a single CPU core (no GPU available in this environment), which made the
project's default finetune budget (1000 epochs, early-stop patience 200,
10 splits per dataset, per config) impractical to run interactively.
Numbers below should be treated as preliminary evidence about training
dynamics, not final benchmark results.

\begin{itemize}
  \item \textbf{Dataset:} Cora, node classification, all 10 splits shipped
        with the dataset (\texttt{finetune.py} always iterates over every
        provided split; there is no independent random seed here).
  \item \textbf{Backbone:} non-MoE GraphSAGE encoder (2 layers,
        hidden dim 768), loaded from the \texttt{default} pretrain
        checkpoint at epoch 50 (\texttt{ckpts/pretrain\_model/default}).
  \item \textbf{Budget:} 60 finetune epochs, early-stop patience 30 --
        far short of the config-file default (1000 epochs / patience
        200). Learning rate and all other hyperparameters use the
        \texttt{config/finetune.yaml} Cora defaults.
  \item \textbf{Coefficients:} a single untuned value, $10^{-3}$, for
        \texttt{decoder\_jac\_coeff} and/or \texttt{encoder\_lip\_coeff}
        in each configuration below. No sweep over coefficient magnitude
        was performed.
  \item \textbf{Metric:} node-classification accuracy (\%), mean $\pm$
        standard deviation across the 10 splits, reported at each split's
        best validation epoch (the project's existing early-stopping
        selection rule).
\end{itemize}

\section{Results}
\label{sec:results}

\begin{center}
\begin{tabular}{lccccc}
\toprule
Configuration & $\lambda_{\text{dec}}$ & $\lambda_{\text{enc}}$ & Train & Val & Test \\
\midrule
Baseline (no regularization)     & 0     & 0     & \score{98.43}{4.25} & \score{76.28}{3.34} & \score{75.87}{2.10} \\
Head-only (prior behavior)       & $10^{-3}$ & 0     & \score{98.43}{4.25} & \score{76.04}{3.35} & \score{75.71}{2.30} \\
Encoder-only (new)               & 0     & $10^{-3}$ & \score{99.64}{0.86} & \score{75.98}{2.82} & \score{75.78}{1.94} \\
Head + Encoder (new, combined)   & $10^{-3}$ & $10^{-3}$ & \score{99.93}{0.21} & \score{75.96}{2.84} & \score{76.03}{0.88} \\
\bottomrule
\end{tabular}
\end{center}

\section{Reading the Results}

\paragraph{Clean-accuracy means are flat.} All four configurations land
within roughly one standard deviation of each other on test accuracy
(75.7\%--76.0\%). At this coefficient and epoch budget, neither penalty
moves mean clean accuracy on Cora in either direction. This is expected:
Lipschitz regularization is a robustness/sensitivity control, not
principally an accuracy lever, and clean i.i.d.\ accuracy is not the metric
that should reveal its effect.

\paragraph{The head-only penalty changes nothing measurable.} The
head-only row is nearly identical to the baseline on every column,
including the across-split standard deviation (train std $4.25 \to 4.25$,
test std $2.10 \to 2.30$). This matches the original concern that motivated
this change: constraining only the decoder's weight norm has little
effect when the unconstrained encoder is free to reshape the input
beforehand.

\paragraph{The encoder-side penalty visibly changes training dynamics.}
Both configurations with $\lambda_{\text{enc}} > 0$ show a sharp drop in
across-split variance on train accuracy (std $4.25 \to 0.86 \to 0.21$) and
a smaller but consistent drop in test-accuracy variance (std
$2.10\text{--}2.30 \to 1.94 \to 0.88$) as the encoder penalty is added and
then combined with the head penalty. In other words, encoder-level
regularization made the 10 finetuning runs converge more consistently with
each other, even though it did not change their average outcome at this
coefficient. That is a real, measured difference in training behavior
attributable specifically to constraining the encoder, since the head-only
run does not reproduce it.

\paragraph{This does not yet test the actual robustness claim.} The
motivation for encoder-level Lipschitz control is bounded sensitivity to
input perturbations (missing features, dropped or shifted edges,
degree/homophily shift) -- not clean-split variance. STEM-GNN already ships
scripts for exactly this (\path{scripts/missing_feature.py},
\path{scripts/random_edge_drop.py},
\path{scripts/degree_shift_ood.py},
\path{scripts/homophily_shift_ood.py}), and none of them were run for
this memo due to the single-CPU compute budget available. The variance
reduction observed here is suggestive but not a substitute for that
evaluation.

\section{Baseline Hyperparameter Optimization (Optuna and Hyperopt)}
\label{sec:optuna}

The comparison in Section~\ref{sec:results} uses the un-tuned
\texttt{config/finetune.yaml} defaults for Cora (learning rate $5\times
10^{-4}$, dropout $0.15$, batch normalization) for every configuration.
Before trusting any accuracy difference between the regularization
settings, it is worth first checking whether the \emph{baseline} model
(no head or encoder regularization, i.e.\ $\lambda_{\text{dec}} =
\lambda_{\text{enc}} = 0$) is even reasonably tuned. We used Optuna to
search the baseline's finetune-stage hyperparameters.

\paragraph{What was tuned.} The pretrained encoder and vector-quantizer
architectures are fixed by the loaded checkpoint (weight shapes must match
the saved \texttt{state\_dict}), so only finetune-stage hyperparameters
that do not change layer shapes are searchable without re-pretraining:
\begin{itemize}
  \item \textbf{Learning rate} (\texttt{finetune\_lr}), log-uniform on
        $[10^{-4}, 10^{-2}]$;
  \item \textbf{Dropout} (\texttt{dropout}), uniform on $[0, 0.5]$;
  \item \textbf{Normalization mode} (\texttt{normalize}), categorical over
        $\{\texttt{none}, \texttt{batch}\}$ (\texttt{layer} was excluded --
        it is accepted by the CLI but the encoder's forward pass applies
        \texttt{BatchNorm1d} regardless of which string is passed, so it is
        not actually a distinct option in the current code).
\end{itemize}
Both regularization coefficients were held fixed at $\lambda_{\text{dec}}
= \lambda_{\text{enc}} = 0$ throughout this search: the goal here is a
better-tuned \emph{baseline}, not a joint search over baseline and
regularization strength.

\paragraph{Implementation.} The search was run twice, once with each
library, over the identical search space and budget described below, so
the two are directly comparable. \path{results/hpo/optuna_baseline.py}
and \path{results/hpo/hyperopt_baseline.py} each build the same parameter
dictionary that \path{finetune.py}'s command-line entry point would build
(Cora defaults from \texttt{config/finetune.yaml}), then import and call
\texttt{finetune.run(params)} directly (in-process, no subprocess spawn)
so that each trial reuses the already-imported PyTorch/PyG modules. Optuna
uses its default TPE sampler and persists every trial to an on-disk SQLite
database (\path{results/hpo/optuna_baseline.db}); Hyperopt uses
\texttt{tpe.suggest} and persists its \texttt{Trials} object to a pickle
file (\path{results/hpo/hyperopt_baseline_trials.pkl}) after every trial.
Both are resumable: if the process is interrupted, re-running the same
script loads the existing state and only runs the remaining trials rather
than starting over. \emph{(This safeguard was added after two earlier
long-running background experiments in this session were silently killed
by the sandbox environment recycling mid-run, losing all unsaved
progress.)}

\paragraph{A caught bug, noted for transparency.} The first Hyperopt run
re-created a fresh, identically-seeded random generator on every
incremental \texttt{fmin} call instead of reusing one generator across
calls. This made every ``next'' suggestion during the startup (random)
phase identical to the first one, so all 12 trials silently evaluated the
exact same hyperparameters instead of exploring the space -- a real bug
in this memo's driver script, not a property of Hyperopt. It was caught
by inspecting the per-trial log before trusting any result (the repeated
parameters and repeated accuracy were the tell), fixed by persisting the
random generator's state alongside the trial history, and re-verified with
a 4-trial smoke test showing genuine variation before the full run below
was produced.

\paragraph{Search budget.} Running the full 10-split, 1000-epoch protocol
per trial is not tractable on the single CPU core available here, so the
search phase uses a reduced budget: a single fixed Cora split
(\texttt{finetune\_seed=0}), 30 finetune epochs, early-stop patience 15,
for 12 trials, optimizing validation accuracy on that split. The
best-found hyperparameters from each library are then re-evaluated once
under the same protocol as Section~\ref{sec:results} (all 10 splits, 60
epochs, early-stop patience 30) to get a number directly comparable to
the baseline row of that table.

\paragraph{Results.} Optuna's 12-trial search found its best single-split
validation accuracy ($73.4\%$) at \texttt{finetune\_lr}$\,\approx
3.47\times 10^{-4}$, \texttt{dropout}$\,\approx 0.014$,
\texttt{normalize}$\,=\,$\texttt{none}. Hyperopt's 12-trial search (same
space and budget, after the bug fix above) found its best single-split
validation accuracy ($75.0\%$) at \texttt{finetune\_lr}$\,\approx
3.56\times 10^{-4}$, \texttt{dropout}$\,\approx 0.094$,
\texttt{normalize}$\,=\,$\texttt{none}. The two libraries agree closely on
learning rate (both land within 3\% of $3.5\times10^{-4}$, versus the
default $5\times10^{-4}$) and on turning batch normalization off, and
disagree more on dropout ($0.014$ vs.\ $0.094$, both well below the
default $0.15$) -- consistent directional signal on two of three
dimensions, with the third under-constrained at this trial budget.
Re-evaluating each configuration under the full 10-split protocol (60
epochs, patience 30) gives:

\begin{center}
\begin{tabular}{lccc}
\toprule
Configuration & Train & Val & Test \\
\midrule
Untuned defaults (Section~\ref{sec:results} baseline) & \score{98.43}{4.25} & \score{76.28}{3.34} & \score{75.87}{2.10} \\
Optuna-tuned baseline                                  & \score{98.50}{4.27} & \score{76.36}{2.85} & \score{75.91}{1.78} \\
Hyperopt-tuned baseline                                & \score{98.50}{4.27} & \score{76.18}{3.10} & \score{75.94}{2.03} \\
\bottomrule
\end{tabular}
\end{center}

\emph{Neither library meaningfully improves on the untuned baseline, and
both land on essentially the same operating point.} All three rows agree
to within $0.2$ points on every column mean -- far inside one standard
deviation -- despite Optuna and Hyperopt proposing visibly different
dropout values ($0.014$ vs.\ $0.094$) to get there. Optuna's run shows the
tightest val/test spread of the three (val std $2.85$, test std $1.78$);
Hyperopt's sits between Optuna's and the untuned default (val std $3.10$,
test std $2.03$). Practically, this says the \texttt{config/finetune.yaml}
Cora defaults were already a reasonable operating point for this
checkpoint and epoch budget -- two independent search algorithms
converged to the same conclusion -- and the head-vs.-encoder
regularization comparison in Section~\ref{sec:results} is not an artifact
of a poorly-tuned baseline.

\section{Test-Time Adaptation and Calibration Under Feature Shift}
\label{sec:tta}

\paragraph{Motivation.} A codebase exploration surfaced this as a natural
complementary mechanism to the Lipschitz penalties above, not a
replacement for them: the penalties in Section~\ref{sec:results} act at
\emph{training} time, bounding how much the encoder and head are allowed
to amplify their input; they say nothing about what happens at
\emph{inference} time on a graph that has already shifted. Two standard,
well-established techniques address exactly that gap without touching
training at all: adapting BatchNorm statistics to the test batch, and
post-hoc temperature scaling.

\paragraph{Implementation.} \texttt{Encoder.encode\_with\_bn\_adaptation}
(\path{model/encoder.py}) runs the forward pass with every
\texttt{BatchNorm1d} layer normalizing against the current batch's own
statistics rather than the stored running statistics from training
(``prediction-time batch normalization,'' Nado et al.\ 2020; Schneider et
al.\ 2020) -- no weights change, and the running statistics are snapshotted
and restored afterward, so this is a pure, reversible per-call adaptation.
\path{utils/calibration.py} adds \texttt{fit\_temperature} (post-hoc
temperature scaling, Guo et al.\ 2017: fit a scalar $T$ on held-out
in-distribution logits, then divide test-time logits by $T$ before
softmax), \texttt{apply\_temperature}, and
\texttt{expected\_calibration\_error} (binned ECE). \texttt{task/node.py}'s
\texttt{eval\_node} gained \texttt{bn\_adapt}, \texttt{temperature}, and
\texttt{return\_ece} keyword arguments (all default off, fully backward
compatible), and \path{finetune.py}'s \texttt{run()} gained an
\texttt{on\_split\_end} callback so an experiment can access the
actually-trained model for extra post-hoc evaluation without duplicating
the training loop.

\paragraph{Experimental setup.} Same reduced budget as the rest of this
memo: the \texttt{default} (non-MoE) checkpoint, 60 finetune epochs,
early-stop patience 30, all 10 Cora splits, $\lambda_{\text{dec}} =
\lambda_{\text{enc}} = 0$ (untuned-baseline architecture, since this
section is about inference-time behavior, not the training-time
penalties). \textbf{One deliberate deviation:} \texttt{normalize} is
forced to \texttt{batch} for this experiment only, overriding Cora's
\texttt{config/finetune.yaml} default of \texttt{none}. Under
\texttt{normalize=none}, \texttt{Encoder.encode} never calls its
\texttt{BatchNorm1d} layers at all, which would make BN adaptation a
no-op by construction -- turning batch normalization on is required to
give it something to adapt. After training, each split's model is
evaluated on: the clean graph; the same graph with a $50\%$ missing-feature
shift applied to validation+test nodes only (\texttt{scripts/}
\texttt{missing\_feature.py}'s \texttt{\_apply\_missing\_features},
\texttt{missing\_prob=0.5}, \texttt{perturb=}\texttt{"valtest"}); and that
shifted graph again under BN adaptation, temperature scaling (fit once on
the \emph{clean} validation logits), and both combined.

\paragraph{Results.} Mean $\pm$ std over the 10 splits, test split only:

\begin{center}
\begin{tabular}{lcc}
\toprule
Condition & Test Accuracy & Test ECE \\
\midrule
Clean (no shift)                     & \score{69.85}{3.74} & \score{0.138}{0.029} \\
Shifted, no adaptation               & \score{62.40}{4.75} & \score{0.125}{0.031} \\
Shifted, BN-adapted                  & \score{66.48}{3.89} & \score{0.132}{0.031} \\
Shifted, temperature-scaled          & \score{62.40}{4.75} & \score{0.090}{0.028} \\
Shifted, BN-adapted + temperature    & \score{66.48}{3.89} & \score{0.096}{0.015} \\
\bottomrule
\end{tabular}
\end{center}

The fitted temperature averaged $T \approx 0.86 \pm 0.10$ across splits.

\paragraph{Reading the results.} The $50\%$ feature shift costs $7.5$
points of test accuracy ($69.85\%\to62.40\%$). \textbf{BN adaptation
recovers about $55\%$ of that drop} ($62.40\%\to66.48\%$, $+4.09$ points)
purely by renormalizing against the shifted batch's own statistics --
no gradient steps, no labels, no change to any weight. \textbf{Temperature
scaling recovers none of the accuracy} ($62.40\%\to62.40\%$, exactly
unchanged, as expected: dividing logits by a scalar is a monotonic
transform and cannot change which class has the highest score) but
\textbf{meaningfully improves calibration} ($0.125\to0.090$ ECE, a $28\%$
relative reduction) even though $T$ was fit only on clean, unshifted
validation data -- the miscalibration direction generalizes across this
shift even though accuracy does not. Combining both gives BN's accuracy
recovery alongside a comparable calibration improvement over the
BN-only condition ($0.132\to0.096$ ECE) and, notably, the tightest spread
of any shifted condition (ECE std $0.015$ vs.\ $0.028$--$0.031$
elsewhere) -- the two adaptations do not appear to interfere with each
other. Two honest caveats: BN adaptation's own ECE is statistically flat
relative to no adaptation ($0.125\to0.132$, within noise) despite the
large accuracy gain, so accuracy recovery and calibration are separate
axes here, not one mechanism producing both; and all of this is
demonstrated only under one shift type (feature dropout), one dataset, and
the same reduced training budget used throughout this memo.

\section{Next Steps}

\begin{enumerate}
  \item Run the existing OOD/perturbation scripts (missing-feature,
        random edge-drop, degree/homophily shift) comparing baseline vs.\
        encoder-only vs.\ combined regularization, which is the setting
        that should actually distinguish the two penalties.
  \item Sweep \texttt{encoder\_lip\_coeff} (e.g.\ $10^{-4}$--$10^{-2}$)
        instead of a single untuned value, since the variance-reduction
        effect observed here may strengthen or trade off against accuracy
        at a different magnitude. (A first attempt at this sweep was lost
        to the same sandbox interruption noted in
        Section~\ref{sec:optuna} and has not yet been re-run.)
  \item Once the Section~\ref{sec:optuna} baseline is tuned, repeat the
        same Optuna search with $\lambda_{\text{enc}}$ included as a
        search dimension, to compare a tuned regularized model against a
        tuned baseline rather than an untuned one.
  \item Repeat on a second dataset (e.g.\ PubMed or WikiCS) and, compute
        permitting, restore the project's standard budget (1000 epochs,
        patience 200) to confirm the trend survives full convergence
        rather than being an artifact of early stopping at 60 epochs.
\end{enumerate}

\end{document}
